\scanpage{32}
\mspage{12}

non plus comprendre les \uline{phénomènes de $p$-torsion cohomologique} des var. alg. en car. $p$. \unsure{D'autres} motivations \inserted{\unsure{pour voir vraiment complet}} ne manquent pas, \deleted{\unsure{peut-être}} \inserted{cf p.ex. l'exposé de Tate sur Brauer, \unsure{pour} heuristique d'une formule pour $\zeta$}

\margin{\uline{par exemple} des groupes de Brauer d'une \unsure{surface} sur un corps fini; \unsure{notre} \unsure{se heurte à} \unsure{un manque de connaissances} de \unsure{cohomologies} ``$p$-adiques'' \unsure{raisonnables}\ldots\ \unsure{ou le global}\ldots}

En analogie avec la coh. $\ell$-adique, on pourrait espérer trouver un foncteur
\diagram{\begin{tikzcd}[ampersand replacement=\&]
X \arrow[r,dash] \& H^i(X)
\end{tikzcd}}
de la catégorie des VA projectives lisses sur $k$, dans celle des modules de t.f. sur $\mathbb Z_p$, ou sur $\mathbb Q_p$, ayant les propriétés habituelles \deleted{\unsure{cf}} (Künneth, dualité, ``bons'' nbs de Betti \ldots). Un argument de Serre montre que déjà pour le $H^1$, et en se bornant aux $X$ \deleted{\unsure{de dim $\leq 2$}}, qui sont des VA de dim $\leq 2$, un tel foncteur n'existe pas. Il semble que l'anneau de coefficients naturel d'une coh. ``$p$-adique'' pour les var. sur un corps $k$ de car $p>0$, est plutôt l'anneau

\margin{31}
